\documentclass[11pt,a4paper]{article}

\usepackage[margin=1.0in]{geometry}
\usepackage{amsmath,amssymb}
\usepackage{booktabs}
\usepackage{tabularx}
\usepackage{hyperref}
\usepackage{microtype}
\usepackage{parskip}
\usepackage{listings}
\usepackage{xcolor}

\lstset{
  language=Python,
  basicstyle=\ttfamily\small,
  keywordstyle=\color{blue!70!black},
  stringstyle=\color{green!50!black},
  commentstyle=\color{gray},
  showstringspaces=false,
  breaklines=true,
  frame=single,
  framesep=4pt,
  xleftmargin=8pt,
  xrightmargin=8pt,
}

\hypersetup{colorlinks=true,linkcolor=blue!60!black,citecolor=blue!60!black,urlcolor=blue!60!black}

\title{From Equivariants to Invariants:\\
       A Primer on SO(3) Representation Theory\\
       for Shape Descriptors}
%\author{}
\date{}

\begin{document}
\maketitle
\tableofcontents
\bigskip

% ---------------------------------------------------------------------------
\section{Introduction}
\label{sec:intro}

Comparing the shapes of 3-D objects requires descriptors that do not depend on
the object's orientation in space.
A \emph{rotation-invariant} descriptor assigns the same number to a shape
regardless of how it is rotated.
Two software packages implement this idea for complementary data types:

\begin{itemize}
  \item \textbf{pykarambola} — computes \emph{Minkowski tensors}
        \cite{SchroderTurk2013,HugSchneider2017}, a family of rank-0, 1, and~2
        tensors derived from a triangulated surface mesh.
  \item \textbf{aics-shparam} — works instead with \emph{spherical harmonic
        coefficients} of a surface parametrisation.
\end{itemize}

Both datasets are \emph{equivariant} under rotation: they do not stay constant,
but they transform in a predictable, structured way.
The central observation of this primer is that both packages follow the same
abstract pipeline for extracting invariants:

\medskip
\centerline{equivariant data
  $\;\longrightarrow\;$ irrep decomposition
  $\;\longrightarrow\;$ contraction to $l=0$
  $\;\longrightarrow\;$ invariant scalar}
\medskip

\paragraph{Terminology stack.}
Four levels appear throughout this primer and are worth naming precisely from the outset.
\begin{enumerate}
  \item \textbf{Equivariants} — the raw input data: Minkowski tensors
        ($w_{000}$, $w_{020}$, \ldots) or spherical harmonic coefficients $f_{l,m}$.
        These transform predictably under rotation but are not themselves invariant.
  \item \textbf{Irreps} — the components obtained by decomposing equivariants into
        irreducible representations (labeled \texttt{'0e'}, \texttt{'2e'}, \ldots).
        Still equivariant, but now in a basis where the rotation action is block-diagonal;
        individual components such as the traceless part of $w_{020}$ carry direct geometric meaning.
  \item \textbf{Invariants} — scalar contractions of irreps that are unchanged under
        rotation, produced by \texttt{compute\_invariants()}.
        Degree-1 invariants (traces of rank-2 tensors) retain clear geometric interpretability;
        degree-2 and degree-3 contractions are algebraically precise but largely opaque geometrically.
  \item \textbf{Generating set vs.\ integrity basis} — a \emph{generating set} is a
        collection of invariants from which every polynomial invariant can be recovered as a
        polynomial in the generators; it may contain redundant members.
        A \emph{minimal integrity basis} is a generating set with all redundancies removed.
        The output of \texttt{compute\_invariants()} is a generating set up to the chosen
        degree, not a minimal basis.
\end{enumerate}

\noindent
Sections~\ref{sec:equivariants}--\ref{sec:parity} develop the mathematics in
increasing detail.
Section~\ref{sec:walkthrough} is a line-by-line walkthrough of the pykarambola
implementation.
Sections~\ref{sec:so2} and~\ref{sec:spharm} cover two extensions:
axis-aware (SO(2)) invariants and spherical-harmonic invariants.

% ---------------------------------------------------------------------------
\section{Equivariants: quantities that transform predictably}
\label{sec:equivariants}

Let $G$ be a group acting on a space $X$ of geometric objects (e.g.\
$G = \mathrm{SO}(3)$, $X = \mathbb{R}^3$).
A function $f : X \to V$ into a vector space $V$ is \emph{equivariant} with
respect to a representation $\rho : G \to \mathrm{GL}(V)$ if
\begin{equation}
  f(g \cdot x) = \rho(g)\, f(x) \qquad \forall\, g \in G,\; x \in X.
  \label{eq:equivariance}
\end{equation}
It is \emph{invariant} when $\rho$ is the trivial representation ($\rho(g) = 1$
for all $g$), i.e.\ $f(g \cdot x) = f(x)$.

\paragraph{Examples at each rank.}
\begin{itemize}
  \item \textbf{Rank-0 (scalar)}: volume $w_{000}$, surface area $w_{100}$, etc.\ are
        invariant (no factor of $R$).
  \item \textbf{Rank-1 (vector)}: the surface centroid $\mathbf{c} = w_{010}/w_{000}$
        satisfies $\mathbf{c}(R\cdot K) = R\,\mathbf{c}(K)$, so it is equivariant
        with $\rho(R) = R$ (the defining representation of SO(3)).
  \item \textbf{Rank-2 (matrix)}: the moment tensor $w_{020}$ satisfies
        $w_{020}(R \cdot K) = R\, w_{020}(K)\, R^\top$, equivariant with
        $\rho(R) = R \otimes R$.
  \item \textbf{Spherical harmonic coefficients}: $f_{l,m}(R \cdot K) =
        \sum_{m'} D^{(l)}_{m'm}(R)\, f_{l,m'}(K)$, equivariant with the
        Wigner $D$-matrix $\rho = D^{(l)}$.
\end{itemize}

\paragraph{Why equivariants suffice.}
Any invariant that is polynomial in the underlying geometric data can be
written as a polynomial contraction of equivariants \cite{Weyl1946}.
So the problem of finding all invariants reduces to: (1) decompose equivariants
into \emph{irreducible representations} (irreps), then (2) contract irreps in
ways that produce an overall $l=0$ output.

% ---------------------------------------------------------------------------
\section{Irreducible Representations of SO(3)}
\label{sec:irreps}

The irreducible representations of SO(3) are labeled by a non-negative integer
$l = 0, 1, 2, \ldots$ called the \emph{angular momentum quantum number}.
An $l$-irrep has dimension $2l+1$; under a rotation $R$ the representation
matrix is the Wigner $D^{(l)}(R)$.

\paragraph{Irrep labels: $l$, \texttt{e}, and \texttt{o}.}
Throughout this document and the pykarambola source code, irreps are labeled by a short string such as \texttt{'0e'}, \texttt{'1e'}, or \texttt{'2e'}.
The integer prefix is the angular momentum number $l$ introduced above.
The letter suffix encodes \emph{parity} under spatial inversion $P : \mathbf{x} \mapsto -\mathbf{x}$, the additional generator that distinguishes $\mathrm{O}(3)$ from $\mathrm{SO}(3)$:
\begin{itemize}
  \item \texttt{e} (\emph{even}, also called \emph{gerade} in the German physics tradition):
        the irrep is unchanged by $P$; parity eigenvalue $+1$.
        True scalars (\texttt{'0e'}), traceless symmetric rank-2 tensors (\texttt{'2e'}), and
        axial vectors (e.g.\ cross products, \texttt{'1e'}) carry even parity.
  \item \texttt{o} (\emph{odd}, also called \emph{ungerade}):
        the irrep acquires a factor of $-1$ under $P$; parity eigenvalue $-1$.
        Polar vectors such as position or displacement are \texttt{'1o'}, and
        pseudo-scalars (e.g.\ the determinant of three polar vectors) are \texttt{'0o'}.
\end{itemize}
The spherical harmonic $Y_l^m$ has parity $(-1)^l$, so irreps built from purely polar-vector inputs naturally carry parity $(-1)^l$: rank-0 and rank-2 tensors are \texttt{e}, rank-1 and rank-3 tensors are \texttt{o}.
Pykarambola departs from this convention by writing \texttt{e} on all tensor labels and handling parity separately through the \texttt{symmetry=} argument; see the label-convention note below.

\paragraph{Rank-2 decomposition.}
A rank-2 symmetric matrix $M$ carries the reducible representation
$\mathbf{1} \otimes \mathbf{1} = \mathbf{0} \oplus \mathbf{2}$ (i.e.\
$l=0 \oplus l=2$).
The decomposition is
\begin{equation}
  M = \underbrace{\frac{\operatorname{Tr} M}{3}\, I}_{l=0,\;\text{isotropic}}
    + \underbrace{\!\left(M_\mathrm{sym} - \frac{\operatorname{Tr} M}{3}\, I\right)\!}_{l=2,\;\text{traceless symmetric}}
  \label{eq:rank2decomp}
\end{equation}
where $M_\mathrm{sym} = (M + M^\top)/2$.
In pykarambola (\texttt{pykarambola/invariants.py}):
\begin{lstlisting}
def _trace_rank2(M):
    return np.trace(M) / 3.0          # l=0 component (float)

def _traceless_rank2(M):
    M_sym = (M + M.T) / 2.0
    return M_sym - (np.trace(M_sym) / 3.0) * np.eye(3)  # l=2 component
\end{lstlisting}

\paragraph{Tensor--irrep table.}
Table~\ref{tab:irreps} shows how each Minkowski tensor is decomposed by
\texttt{decompose\_all()} and what code keys it produces.

\begin{table}[h]
\centering
\caption{Minkowski tensor irrep decomposition.
         \texttt{decompose\_all()} returns entries keyed by
         \texttt{(name, irrep\_label)}.}
\label{tab:irreps}
\begin{tabularx}{\linewidth}{llllX}
\toprule
Tensor & Rank & Physical meaning & Irreps & Code keys \\
\midrule
\texttt{w000} & 0 & Volume & $l=0$ & \texttt{('w000','0e')} \\
\texttt{w010} & 1 & Volume moment (centroid $\times$ vol) & $l=1$ & \texttt{('w010','1e')} \\
\texttt{w020} & 2 & Moment tensor & $l=0 \oplus l=2$ & \texttt{('w020','0e')}, \texttt{('w020','2e')} \\
\texttt{w102} & 2 & Normal distribution & $l=0 \oplus l=2$ & \texttt{('w102','0e')}, \texttt{('w102','2e')} \\
\bottomrule
\end{tabularx}
\end{table}

\paragraph{Label convention.}
pykarambola labels all $l=1$ irreps as \texttt{'1e'}, regardless of parity
(polar vs.\ axial vectors).
This differs from the e3nn library \cite{Geiger2022e3nn}, which uses
\texttt{'1o'} for polar (odd-parity) and \texttt{'1e'} for axial (even-parity)
vectors.
In pykarambola, parity is controlled by the \texttt{symmetry=} argument to
\texttt{compute\_invariants()}, not by the irrep label.

\paragraph{Spherical harmonic coefficients.}
The $f_{l,m}$ produced by aics-shparam are already irreps by construction —
each coefficient band $l$ transforms as a $(2l+1)$-dimensional $D^{(l)}$
representation.
No further decomposition step is required.

% ---------------------------------------------------------------------------
\section{The Contraction Principle: coupling irreps to scalars}
\label{sec:contraction}

\paragraph{Clebsch--Gordan selection rule.}
When two irreps of angular momenta $l_1$ and $l_2$ are coupled, the possible
outputs are $|l_1 - l_2| \le L \le l_1 + l_2$.
To obtain a scalar ($L=0$) we need $l_1 = l_2$ (and a matching parity).
This immediately constrains which contractions are possible:

\medskip
\begin{center}
\renewcommand{\arraystretch}{1.3}
\begin{tabular}{cll}
\toprule
Coupling & Contraction & pykarambola function \\
\midrule
\multicolumn{3}{l}{\textit{Degree~2 (bilinear)}} \\
$1 \otimes 1 \to 0$ & dot product $\mathbf{u} \cdot \mathbf{v}$ & \texttt{\_vector\_dot\_products} \\
$2 \otimes 2 \to 0$ & Frobenius product $\operatorname{Tr}(T^\top S)$ & \texttt{\_frobenius\_inner\_products} \\
\multicolumn{3}{l}{\textit{Degree~3 — O(3) true scalars}} \\
$1 \otimes 2 \otimes 1 \to 0$ & quadratic form $\mathbf{u}^\top T\, \mathbf{v}$ & \texttt{\_quadratic\_forms} \\
$2 \otimes 2 \otimes 2 \to 0$ & triple trace $\operatorname{Tr}(TUS)$ & \texttt{\_triple\_traces} \\
\multicolumn{3}{l}{\textit{Degree~3 — SO(3) pseudo-scalars only}} \\
$1 \otimes 1 \otimes 1 \to \tilde{0}$ & determinant $\det[\mathbf{u},\mathbf{v},\mathbf{w}]$ & \texttt{\_triple\_vector\_dets} \\
$2 \otimes 2 \otimes 1 \to \tilde{0}$ & commutator $[T,U]_{\text{axial}} \cdot \mathbf{v}$ & \texttt{\_commutator\_pseudoscalars} \\
\bottomrule
\end{tabular}
\end{center}
\medskip

\paragraph{Why each contraction gives $l=0$.}
For degree-2 couplings the argument is direct: the Frobenius product
$\operatorname{Tr}(T^\top S)$ satisfies
\[
  \operatorname{Tr}\!\bigl((RT R^\top)^\top (RS R^\top)\bigr)
  = \operatorname{Tr}\!\bigl(R T^\top R^\top R S R^\top\bigr)
  = \operatorname{Tr}(T^\top S)
\]
using $R^\top R = I$ and the cyclic property of trace.
The same argument applies to dot products and to all degree-3 O(3) contractions.

\paragraph{Commutator pseudo-scalar.}
For two traceless symmetric matrices $T_a$, $T_b$, the commutator
$C = T_a T_b - T_b T_a$ is antisymmetric.
Every antisymmetric $3\times3$ matrix has an associated axial vector
$\mathbf{a} = (C_{12}, C_{20}, C_{01})$ such that $C_{ij} = \varepsilon_{ijk} a_k$.
The dot product $\mathbf{a} \cdot \mathbf{v}$ is invariant under proper rotations
but changes sign under a reflection, making it a pseudo-scalar.

\paragraph{Degree-1 invariants.}
The $l=0$ component of any tensor is already an invariant (it is a scalar that
transforms trivially).
For rank-0 tensors this is the value itself; for rank-2 tensors it is the trace
$\operatorname{Tr}(M)/3$.
One linear dependency is removed: $\operatorname{Tr}(w_{102})/3 = w_{100}$
and $\operatorname{Tr}(w_{202})/3 = w_{200}$ hold for any closed mesh
(since $|\hat{n}|^2 = 1$).
When \texttt{deduplicate\_scalars=True} (the default), \texttt{compute\_invariants}
drops these two redundant traces.

% ---------------------------------------------------------------------------
\section{Parity: O(3) versus SO(3)}
\label{sec:parity}

The orthogonal group $\mathrm{O}(3)$ includes reflections in addition to
rotations.
A function invariant under $\mathrm{O}(3)$ must also be unchanged by any
reflection.
\emph{Pseudo-scalars} (the determinant and commutator contractions above) are
invariant under proper rotations but change sign under an improper rotation
(reflection or roto-reflection).
They therefore appear in SO(3)-invariant feature sets but not O(3)-invariant ones.

\begin{center}
\begin{tabular}{ll}
\toprule
Feature subset & Symmetry \\
\midrule
\texttt{symmetry='O3'} & True scalars only (no \texttt{det\_}, no \texttt{comm\_}) \\
\texttt{symmetry='SO3'} & True scalars $\cup$ pseudo-scalars \\
\bottomrule
\end{tabular}
\end{center}

For the standard 14-tensor set (4 rank-0, 4 rank-1, 6 rank-2 tensors) the
difference amounts to $64$ features: $4$ triple-vector determinants and
$60$ commutator pseudo-scalars, giving $155$ features for O(3) and $219$ for
SO(3).

% ---------------------------------------------------------------------------
\section{pykarambola Implementation Walkthrough}
\label{sec:walkthrough}

All invariant code lives in \texttt{pykarambola/invariants.py}.
The public entry point is \texttt{compute\_invariants()}.

\subsection*{Step 1 — Decompose into irreps}

\begin{lstlisting}
def decompose_all(tensors_dict):
    result = {}
    for name, tensor in tensors_dict.items():
        rank = _infer_rank(tensor)
        if rank == 0:
            result[(name, '0e')] = float(tensor)
        elif rank == 1:
            result[(name, '1e')] = np.asarray(tensor, dtype=float)
        elif rank == 2:
            M = np.asarray(tensor, dtype=float)
            result[(name, '0e')] = _trace_rank2(M)    # Tr(M)/3
            result[(name, '2e')] = _traceless_rank2(M)
    return result
\end{lstlisting}

The function returns a dictionary keyed by \texttt{(tensor\_name, irrep\_label)}.
A rank-2 tensor produces two entries; a rank-0 or rank-1 tensor produces one.

\subsection*{Step 2 — Contract by degree}

\texttt{compute\_invariants()} calls the contraction helpers in order:

\begin{lstlisting}
def compute_invariants(tensors_dict, max_degree=3,
                       symmetry='SO3', deduplicate_scalars=True):
    decomposed = decompose_all(tensors_dict)
    result = {}
    result.update(_degree1_scalars(decomposed, deduplicate=deduplicate_scalars))
    if max_degree >= 2:
        result.update(_vector_dot_products(decomposed))
        result.update(_frobenius_inner_products(decomposed))
    if max_degree >= 3:
        result.update(_quadratic_forms(decomposed))
        result.update(_triple_traces(decomposed))
        if symmetry == 'SO3':
            result.update(_triple_vector_dets(decomposed))
            result.update(_commutator_pseudoscalars(decomposed))
    return result
\end{lstlisting}

\subsection*{Worked example}

Given the small tensor set
\begin{lstlisting}
tensors = {
    'w000': 1.0,                         # rank 0
    'w010': np.array([0.5, 0.5, 0.5]),   # rank 1
    'w020': np.eye(3) * 0.1,             # rank 2
}
inv = compute_invariants(tensors, max_degree=2, symmetry='O3')
\end{lstlisting}
the pipeline produces:

\begin{center}
\begin{tabular}{lll}
\toprule
Key & Source & Formula \\
\midrule
\texttt{w000} & degree-1, rank-0 & $w_{000} = 1.0$ \\
\texttt{w020} & degree-1, rank-2 trace & $\operatorname{Tr}(w_{020})/3 = 0.1$ \\
\texttt{dot\_w010\_w010} & degree-2 & $\mathbf{w}_{010} \cdot \mathbf{w}_{010} = 0.75$ \\
\texttt{frob\_w020\_w020} & degree-2 & $\operatorname{Tr}(T_{\text{tl}}^\top T_{\text{tl}})$
  (zero here, since \texttt{w020} is a multiple of $I$) \\
\bottomrule
\end{tabular}
\end{center}

\noindent
Because $w_{020} = 0.1\,I$, its traceless part $T_{\text{tl}}$ is zero and the
Frobenius product is~0.

\subsection*{Feature count for the standard 14-tensor set}

With all 14 standard Minkowski tensors (4 rank-0, 4 rank-1, 6 rank-2),
\texttt{deduplicate\_scalars=True}, the feature counts below are verified by
\texttt{tests/test\_invariants.py::TestInvariantCounts}.

\begin{table}[h]
\centering
\caption{Invariant feature counts for the 14-tensor standard set.
         Columns: degree-1 = scalars/traces; degree-2 = dot products and Frobenius
         products; degree-3 = all cubic contractions.}
\label{tab:counts}
\begin{tabular}{lrrrr}
\toprule
Symmetry & Degree 1 & Degree 2 & Degree 3 & Total \\
\midrule
O(3)  & 8 & 31 & 116 & 155 \\
SO(3) & 8 & 31 & 180 & 219 \\
\bottomrule
\end{tabular}
\end{table}

\noindent
\textbf{Degree-1} (8): four rank-0 scalars plus six matrix traces, minus two
removed by deduplication ($\operatorname{Tr}(w_{102})/3 = w_{100}$ and
$\operatorname{Tr}(w_{202})/3 = w_{200}$).

\textbf{Degree-2} (31): $\binom{4}{2}+4 = 10$ vector dot products plus
$\binom{6}{2}+6 = 21$ Frobenius products.

\textbf{Degree-3, O(3)} (116): $6 \times 10 = 60$ quadratic forms (6 matrices
$\times$ 10 unordered vector pairs) plus $\binom{8}{3} = 56$ triple traces
($\binom{n+r-1}{r}$ with $n=6$ matrices, $r=3$).

\textbf{Degree-3, SO(3) additional} (64): $\binom{4}{3} = 4$ triple-vector
determinants plus $\binom{6}{2} \times 4 = 60$ commutator pseudo-scalars.

% ---------------------------------------------------------------------------
\section{Extending to SO(2): axis-aware invariants}
\label{sec:so2}

When the $z$-axis carries physical meaning — cell polarity, optical axis,
gravitational direction — the relevant symmetry group shrinks from SO(3)
to SO(2) (rotations about $z$).
This breaks each SO(3) irrep into \emph{charges} $m$ under $z$-rotation:
a rank-2 tensor contributes $m=0$, $|m|=1$, and $|m|=2$ components.

\texttt{\_decompose\_so2()} produces the following labels per tensor:

\begin{center}
\begin{tabular}{lll}
\toprule
Label & Content & Rank \\
\midrule
\texttt{'sc'} & scalar value & 0 \\
\texttt{'z'} & $v_z$ & 1 \\
\texttt{'xy'} & $[v_x, v_y]$ doublet ($|m|=1$) & 1 \\
\texttt{'tr'} & $\operatorname{Tr}(M)/3$ isotropic & 2 \\
\texttt{'tzz'} & $M_{zz} - \operatorname{Tr}(M)/3$ ($m=0$ traceless) & 2 \\
\texttt{'xz'} & $[M_{xz}, M_{yz}]$ ($|m|=1$) & 2 \\
\texttt{'m2'} & $[M_{xx}-M_{yy},\, 2M_{xy}]$ ($|m|=2$) & 2 \\
\bottomrule
\end{tabular}
\end{center}

SO(2)-invariant contractions work by charge conservation:
\begin{itemize}
  \item \textbf{Degree-2}: dot products and wedge products of same-charge doublets
        (\texttt{d1\_}, \texttt{x1\_} for $|m|=1$; \texttt{d2\_}, \texttt{x2\_} for $|m|=2$).
        For doublets $\mathbf{a} = a_x + i\,a_y$ and $\mathbf{b} = b_x + i\,b_y$
        (read as complex numbers), $\operatorname{Re}(\bar{a}b) = \mathbf{a}\cdot\mathbf{b}$
        and $\operatorname{Im}(\bar{a}b) = a_x b_y - a_y b_x$ are both SO(2)-invariant.
  \item \textbf{Degree-3}: $|m|=1 \otimes |m|=1 \otimes |m|=2$ couplings via complex
        multiplication (\texttt{tp\_re\_}, \texttt{tp\_im\_}).
        For $\mathbf{a},\mathbf{b}$ ($|m|=1$) and $\mathbf{c}$ ($|m|=2$):
        $\operatorname{Re}(\bar{\mathbf{c}} \cdot \mathbf{a}\mathbf{b})$ and
        $\operatorname{Im}(\bar{\mathbf{c}} \cdot \mathbf{a}\mathbf{b})$, where
        $\mathbf{a}\mathbf{b}$ is complex multiplication.
\end{itemize}

For the standard 14-tensor set the total SO(2) feature count is $18 + 136 + 660 = \mathbf{814}$,
verified by \texttt{TestSO2Invariants.test\_so2\_invariant\_count\_14\_tensors}.
O(2) is obtained by filtering out features that are odd under $z \to -z$
(the $v_z$ degree-1 components, mixed-parity $|m|=1$ bilinears, and
mixed-parity triple products).

% ---------------------------------------------------------------------------
\section{Parallel: Spherical Harmonic Invariants}
\label{sec:spharm}

The aics-shparam package parametrises a surface as a spherical harmonic expansion
$f(\theta,\phi) = \sum_{l,m} f_{l,m}\, Y_l^m(\theta,\phi)$ and returns the
coefficient array $\{f_{l,m}\}$.
Because $Y_l^m$ itself transforms as a $D^{(l)}$ irrep, the coefficients are
already decomposed into SO(3) irreps — no further decomposition step is needed.
Each band $l$ is a $(2l+1)$-dimensional irrep; the table below shows what these
look like for the first few values of $l$.

\begin{center}
\begin{tabular}{cll}
\toprule
$l$ & Dimension ($2l+1$) & Geometric character \\
\midrule
0 & 1 & scalar \\
1 & 3 & 3-vector \\
2 & 5 & traceless symmetric $3\times3$ matrix (6 entries minus 1 trace constraint) \\
3 & 7 & rank-3 traceless symmetric object \\
\bottomrule
\end{tabular}
\end{center}

\noindent
Only $l=1$ coincides with an ordinary 3-vector.
For $l \ge 2$ the irrep is stored as a higher-rank array but has fewer independent
components than its ambient container: the $l=2$ irrep fits in a $3\times3$ matrix
yet has only 5 degrees of freedom.

The same contraction principle then applies directly:

\begin{itemize}
  \item \textbf{Power spectrum} (degree-2): $S_l = \sum_{m=-l}^{l} |f_{l,m}|^2$
        is the $l \otimes l \to 0$ Frobenius contraction; it is invariant under
        SO(3).
  \item \textbf{Bispectrum} (degree-3): $B_{l_1 l_2 l} =
        \operatorname{Re}\!\sum_{m_1, m_2} C^{l,m}_{l_1,m_1;\, l_2,m_2}\,
        f_{l_1,m_1} f_{l_2,m_2} f^*_{l,m}$ with $m = m_1 + m_2$, where
        $C^{l,m}_{l_1,m_1;\,l_2,m_2}$ are Clebsch--Gordan coefficients.
        This is a degree-3 contraction coupling three irreps to $L=0$.
\end{itemize}

The structure is identical to Sections~\ref{sec:contraction}--\ref{sec:parity}:
irrep channels, degree-$k$ couplings via Clebsch--Gordan coefficients, and the
same parity distinction between SO(3) and O(3) invariants.

% ---------------------------------------------------------------------------
\section{Key Takeaways}
\label{sec:takeaways}

\begin{enumerate}
  \item \textbf{Equivariants are enough.}  Any polynomial invariant can be
        written as a contraction of equivariants; the pipeline decompose $\to$
        contract is systematic and (up to chosen degree) exhaustive
        \cite{KondorTrivedi2018}.

  \item \textbf{Degree controls complexity.}
        Degree-1 features are free (linear in the input tensors);
        degree-2 features are bilinear (products of two irreps);
        degree-3 features are cubic.
        Higher-degree contractions exist but grow combinatorially.

  \item \textbf{Generating set, not a minimal basis.}
        The raw Minkowski tensors are equivariants, not invariants.
        The output of \texttt{compute\_invariants()} is a generating set of invariants:
        it contains all polynomial invariants up to the chosen degree, but it may include
        algebraic dependencies (e.g.\ Cayley--Hamilton relations among triple traces).
        Apart from the two trace identities removed by \texttt{deduplicate\_scalars},
        redundancies are not eliminated.
        A minimal integrity basis — the smallest generating set — is known to exist by the
        Hilbert basis theorem but is not computed.

  \item \textbf{Pseudo-scalars require a parity decision.}
        The determinant and commutator contractions are invariant under
        $\mathrm{SO}(3)$ but not $\mathrm{O}(3)$.
        Use \texttt{symmetry='O3'} to exclude them when reflective symmetry is
        physically meaningful.

  \item \textbf{Both packages share the abstract structure.}
        Whether the input is Minkowski tensors or spherical harmonic coefficients,
        the invariant construction is: irrep decomposition $\to$ degree-$k$
        Clebsch--Gordan coupling $\to$ scalar.
\end{enumerate}

\bibliographystyle{plain}
\bibliography{invariants_primer}

\end{document}
